\documentclass[11pt]{article}
\usepackage[paperwidth=210mm,paperheight=5000mm,margin=20mm]{geometry}
\usepackage{amsmath,amssymb,xcolor,graphicx}
\usepackage[hypcap=false]{caption}
\usepackage[hidelinks,bookmarksnumbered,bookmarksopen]{hyperref}
\hypersetup{pdftitle={Features \& tests: 100\% of \$x\$},pdfauthor={Dr Ø. Ñandú},
  pdfcreator={KherveNote}}
\renewcommand\labelitemii{\textbullet}
\renewcommand\labelitemiii{\textbullet}
\renewcommand\labelitemiv{\textbullet}
\renewcommand\labelenumii{\arabic{enumi}.\arabic{enumii}.}
\renewcommand\labelenumiii{\arabic{enumi}.\arabic{enumii}.\arabic{enumiii}.}
\renewcommand\labelenumiv{\arabic{enumi}.\arabic{enumii}.\arabic{enumiii}.\arabic{enumiv}.}
\setlength{\parindent}{0pt}
\setlength{\parskip}{0.6em}
\definecolor{knoteaccent}{HTML}{1A6DD8}
\definecolor{knotemuted}{HTML}{6B6B6B}
\definecolor{knotekey}{HTML}{D96B00}
\newcommand\knotetime[1]{\leavevmode\llap{\color{knotemuted}\scriptsize #1\hspace{1em}}}
\newcommand\knotekey[2][]{\par\noindent#1\colorbox{knotekey!12}{\parbox{\dimexpr\linewidth-2\fboxsep}{%
  \textbf{\color{knotekey}$\star$ Key point.} #2}}\par}
\newcommand\knotequestion[2][]{\par\noindent#1{\color{knoteaccent}\textbf{?}}~\emph{#2}\par}
\newenvironment{knotetranscript}{\par\color{knotemuted}\small}{\par}
\newcommand\knotesummary[1]{\par\noindent\fcolorbox{knoteaccent}{knoteaccent!6}{%
  \parbox{\dimexpr\linewidth-2\fboxsep-2\fboxrule}{\textbf{Summary}\par #1}}\par}
\makeatletter
\newdimen\knote@limit \knote@limit=4950mm
\newbox\knote@page \newbox\knote@chunk
\def\knote@ht#1{\dimexpr\ht#1+\dp#1\relax}
\def\knote@setheight#1{\global\pdfpageheight=#1\relax
  \special{pdf:pagesize width \the\paperwidth\space height \the\pdfpageheight}}
\def\knote@flush{\ifvoid\knote@page\else
  \knote@setheight{\dimexpr\knote@ht\knote@page+40mm+2mm\relax}%
  \box\knote@page\par\clearpage\global\pdfpageheight=\paperheight\fi}
\newenvironment{knotechunk}
  {\par\setbox\knote@chunk=\vbox\bgroup\hsize=\textwidth\linewidth=\textwidth}
  {\par\egroup
   \ifdim\knote@ht\knote@chunk>\knote@limit
     \knote@flush\unvbox\knote@chunk\par\clearpage
   \else\ifvoid\knote@page
     \global\setbox\knote@page=\vbox{\unvbox\knote@chunk}%
   \else\ifdim\dimexpr\knote@ht\knote@page+\knote@ht\knote@chunk\relax>\knote@limit
     \knote@flush\global\setbox\knote@page=\vbox{\unvbox\knote@chunk}%
   \else
     \global\setbox\knote@page=\vbox{\unvbox\knote@page\unvbox\knote@chunk}%
   \fi\fi\fi}
\AtEndDocument{\knote@flush}
\makeatother
\pagestyle{empty}
\begin{document}
\begin{knotechunk}
{\color{knoteaccent}\rule{\linewidth}{1.5pt}}\par
{\LARGE\bfseries Features \& tests: 100\% of $x$}\par
{\large Dr Ø. Ñandú}\hfill{\color{knotemuted}7 October 2026 \textperiodcentered{} Room \#4 \{B\}}\par
{\color{knoteaccent}\rule{\linewidth}{0.6pt}}\par
\knotesummary{\begin{itemize}\setlength{\itemsep}{0pt}\setlength{\parskip}{0pt}
  \item first point
  \item second with $E=mc^2$
\end{itemize}

A closing line.}

Lead-in before any section, at 12:30.
\end{knotechunk}

\begin{knotechunk}
\section{Marks \& maths}

\underline{\textbf{Bold}}\underline{, }\underline{\textit{italic}} and \underline{underlined} words, then line\newline
breaks.

It costs \$5 and \$10; inline $E = h\nu - \phi$ and $k_1$ and \[k = A\exp(-E_a/RT)\] then \[\int_0^1 x\,dx\] end.

Symbols: \ensuremath{\alpha} \ensuremath{\beta} \ensuremath{\Delta} \ensuremath{\approx} \ensuremath{\rightarrow} ° × ± µ H\textsubscript{2}O Al\textsuperscript{3+} cm\textsuperscript{-1} \textasciitilde{} \textasciicircum{} \textbackslash{} \_ \# \& \% \{ \} and $\lambda  = 2^{\circ}$.

\knotekey{Remember: $\Delta G < 0$.}

\knotequestion{Why is it **so**?}

\begin{knotetranscript}
\knotetime{01:30}so um the next slide shows
\end{knotetranscript}

\subsection{A subsection}

\subsubsection{A sub-subsection}
\end{knotechunk}

\begin{knotechunk}
\section{Lists}

\begin{enumerate}\setlength{\itemsep}{0pt}\setlength{\parskip}{0pt}
  \item first
  \begin{enumerate}\setlength{\itemsep}{0pt}\setlength{\parskip}{0pt}
    \item nested
    \begin{itemize}\setlength{\itemsep}{0pt}\setlength{\parskip}{0pt}
      \item deeper bullet
    \end{itemize}
    \item \textbf{back} to two
  \end{enumerate}
  \item second
\end{enumerate}
\begin{itemize}\setlength{\itemsep}{0pt}\setlength{\parskip}{0pt}
  \item a bullet list now
\end{itemize}

Typed list:
\begin{itemize}\setlength{\itemsep}{0pt}\setlength{\parskip}{0pt}
  \item one
  \begin{itemize}\setlength{\itemsep}{0pt}\setlength{\parskip}{0pt}
    \item one.a
  \end{itemize}
\end{itemize}
\begin{enumerate}\setlength{\itemsep}{0pt}\setlength{\parskip}{0pt}
  \item numbered
\end{enumerate}

New paragraph.

\begin{center}
\includegraphics[width=0.85\linewidth,height=110mm,keepaspectratio]{assets/abc123.png}
\captionof{figure}{A caption $x^2$}
\end{center}

{\color{knotemuted}\textbf{Attached document:} Slides.pdf}
\end{knotechunk}
\end{document}
